% =============================================================================
% 文件: sections/s8_evaluation.tex   第八章 模型评价与推广
% =============================================================================
\section{模型评价与推广}

\subsection{模型优点}

\begin{enumerate}[leftmargin=2.2em,itemsep=0.3em]
  \item \textbf{机理清晰、四问统一。}四个问题共用同一套控制方程与同一个求解器，
        差异仅体现在物性关系、几何与输出配置上。这种“一个模型多组配置”的结构
        使得任何一处改进都自动作用于全部问题，便于交叉验证，也避免了不同问题
        各写一套代码时常见的口径不一致。
  \item \textbf{数值格式与物理性质相容。}界面扩散系数取调和平均、容量系数与解
        与隐式 Euler 的正容量、对角占优结构共同保证基准算例的离散稳定性；
        含水率非负且最大值单调下降。温度相对于附件 1 的全局最高温度未越界，
        但平台边界切换后表面重构值可能有很小的瞬态过冲。水量收支的相对闭合误差
        小于 $10^{-11}$，说明有限体积通量在离散意义下守恒。
  \item \textbf{边界重构提高了表层精度。}用半单元“扩散阻力—表面膜阻力串联”
        重构真实表面值，解决了“最后一个单元中心不等于表面”的常见错误；
        与解析解对照表明，重构后表面误差显著小于直接取最外单元中心值；当前
        半单元重构的边界误差仍呈 $O(\Delta r)$，因此不能宣称与内部点严格同阶。
  \item \textbf{误差可追溯、结论有条件。}论文明确给出网格数、时间步、
        边界处理与终点判据，并用解析解、两种时间积分格式、守恒律、
        合成数据回归四条独立证据链验证；干燥时长以区间加主导不确定源的形式
        给出，而不是一个看似精确的单点数字。
  \item \textbf{对模型假设本身做了检验。}附件 2 的收缩轨迹与附录 4 的密度经验式
        之间的不自洽被定量地揭示出来，使问题 4 的结论被正确地限定为
        “给定收缩轨迹下的预测”，避免了对模型适用范围的外推。
\end{enumerate}

\subsection{模型缺点与改进方向}

\begin{enumerate}[leftmargin=2.2em,itemsep=0.3em]
  \item \textbf{传质边界驱动势的单位口径未经实验标定。}
        题目给出的 $h_m=8\times10^{-7}\ \mathrm{m/s}$ 与“烘房水分浓度
        （kg/kg）”在量纲上需要一个隐含的换算约定。本文按“$h_m$ 已吸收单位换算”
        处理，并用 $C_{\mathrm{air}}$ 的整体缩放与 $h_m$ 的合成反演量化了该假设的
        影响（见 7.5 节）。
        \textbf{改进方向}：若能获得药材的吸附等温线，应把边界条件改写为
        $-D\partial C/\partial r=h_m^{\ast}(\rho_{v,s}-\rho_{v,\infty})$，
        其中 $\rho_v$ 为水蒸气质量浓度，并由表面温度查饱和湿度确定平衡含水率。
  \item \textbf{未计入汽化潜热与内部蒸发。}基准模型把热边界简化为导热加对流。
        代码保留可选的 $+\rho_{\mathrm{d}}L_v\partial C/\partial t$ 诊断项，
        但由于题面未给出相应的气相平衡与质量通量标定，它不纳入基准结论。
        \textbf{改进方向}：引入“湿球温度”形式的耦合边界，把表面蒸发与
        表面能量平衡联立，并利用物料表面温度的实验数据标定。
  \item \textbf{收缩与密度经验式不自洽。}附件 2 的实测半径轨迹与附录 4 的
        $\rho(C)$ 无法由同一个固体质量守恒同时满足（偏差达百分之十几）。
        \textbf{改进方向}：把 $R(t)$ 作为待求量，增加固相质量守恒与力学本构
        （例如把收缩率写成含水率的函数并标定），从而把“轨迹输入”变成“轨迹输出”。
  \item \textbf{一维径向假设忽略了端面效应。}端面面积占比约 $8\%$，
        轴向湿度梯度在干燥后期可能不可忽略。
        \textbf{改进方向}：推广为二维轴对称 $(r,z)$ 模型，或对端面单独加一个
        “等效长度”修正，使总传质面积与真实几何一致。
  \item \textbf{物性经验式的外推风险。}附录 3、附录 4 的 $\mathrm{e}^{-a/C}$
        在 $C\to0$ 时趋于零；终点判据取 $C=0.15$ 位于公式声明区间边界，
        更低含水率处仍属于外推。
        \textbf{改进方向}：用分段或饱和型函数（例如
        $D$ 的低含水率分段延拓替换简单指数型，并用等温吸附实验标定。
\end{enumerate}

\subsection{模型的推广}

本文的模型框架并不限于中药材烘干，可直接推广到如下场景：

\begin{enumerate}[leftmargin=2.2em,itemsep=0.25em]
  \item \textbf{其他长径比较大的圆柱形物料的干燥}：如中药材根茎切片、
        粮食颗粒的柱状近似、木材圆柱试件、电缆绝缘层的水分扩散等，
        只需替换物性经验关系与环境边界。
  \item \textbf{带收缩的物料干燥}：物质坐标形式对“收缩由均匀失水引起”的
        物料（多数果蔬、中药材）都适用，只需给出 $R(t)$ 或给出
        $R$ 与含水率的关系式。
  \item \textbf{瞬态热传导/扩散的通用求解器}：把物性函数替换为常数，
        本求解器即退化为标准圆柱导热/扩散求解器；
        把 $R(t)$ 设为常数即为定域问题，把 $\theta$ 设为 $1/2$ 即得二阶时间精度。
  \item \textbf{工艺优化与控制器设计}：模型可作为“数字孪生”内核，
        把烘房温度、湿度作为控制变量，以“干燥时间最短、能耗最低、
        品质指标损失最小”为目标做多目标优化，得到分段变温变湿的干燥曲线。
\end{enumerate}
